En esta sección se explicará el camino desde los archivos de audio hasta los \textit{batches} de entrenamiento.

\begin{itemize}
    \item \textbf{1. Organización inicial del dataset:}
    El dataset se compone de mezclas y fuentes separadas. Para cada pista se tiene:
    \begin{itemize}
        \item \textit{mixture.wav}: la mezcla original.
        \item \textit{vocals.wav}: pista para las vocales.
        \item \textit{drums.wav}: pista para la batería.
        \item \textit{bass.wav}: pista para el bajo.
        \item \textit{others.wav}: pista para elementos extra. 
    \end{itemize}

    \item \textbf{2. Construcción de escenarios:}
    Como ya se ha comentado en secciones anteriores, el modelo contempla dos escenarios de salida: cuatro y dos pistas. Partiendo de ello, el dataset de entrada para los dos tipos de entrenamiento debe ser acorde a la salida. Para ello:

    \begin{itemize}
        \item \textit{2 stems}: se mantienen las vocales y se agrupa el resto como acompañamiento.
        \item \textit{4 stems}: se usan las cuatro fuentes originales.
    \end{itemize}

    \item \textbf{3. Conversión a dominio espectral:}
    Se calcula la STFT de mezcla y fuentes con los mismos parámetros. Se extraen magnitudes para entrenamiento y se conserva la fase de la mezcla para reconstrucción/evaluación.

    \item \textbf{4. Preparación de \textit{batches}:}
    Cada \textit{batch} contiene al menos:
    \begin{itemize}
        \item \texttt{mix\_magnitude}
        \item \texttt{source\_magnitudes}
    \end{itemize}
    
    Y para pérdidas que piden fase:
    \begin{itemize}
        \item \texttt{mixture\_phase}
        \item \texttt{source\_phase}
    \end{itemize}

    La versión final del proceso de entrenamiento comprueba la presencia de estos dos últimos para las funciones que las piden, y usa \texttt{source\_magnitudes} como \textit{target} principal durante el entrenamiento.

    \item \textbf{5. División train/validation/test:}
    \begin{itemize}
        \item \textit{train/validation}: usado durante entrenamiento y \textit{early stopping}.
        \item \texttt{representative\_test}: conjunto fijo usado solo para evaluación final.
    \end{itemize}
    División 60-30.

    \item \textbf{6. Normalización/canales:}
    El sistema se configuró para un canal de audio, lo que reduce el coste computacional y se recomienda en la documentación del modelo. También simplifica el proceso de comparación entre fuentes.

\end{itemize}